\begin{center}
SGA 7

\vspace{1em}
\underline{Exposé IX}

\vspace{1em}
\underline{\underline{MODELOS DE NÉRON Y MONODROMÍA}}

\vspace{1em}
\underline{por A. Grothendieck}

(con un apéndice de M. Raynaud)
\end{center}

\subsection*{Índice}

\begin{enumerate}[label=\arabic*.]
\setcounter{enumi}{-1}
\item Introducción
\item Modelo de Néron de los esquemas abelianos: notación; el
  emparejamiento canónico
  $\Phi_o\times\Phi'_o\longrightarrow(\mathbb Q/\mathbb Z)_k$, y la
  biextensión canónica $W^o$ de $(A^o,{A'}^o)$
\item Parte fija y parte tórica de $T_\ell(A_K)$. Criterios de buena
  reducción. Teorema de ortogonalidad para $\ell\ne p$
\item El caso de reducción semiestable. El teorema de reducción
  semiestable
\item Aplicación a una conjetura de Serre--Tate y al conductor
\item El teorema de ortogonalidad para $\ell=p$ (caso semiestable).
  Dualidad de los esquemas abelianos $B_o$ y $B'_o$. Caracterización de
  la parte fija. Criterios de buena reducción, reducción esencialmente
  buena y reducción semiestable
\item Observaciones sobre la construcción de la parte tórica y la parte
  fija de $T_p(A)$ en el caso de reducción no semiestable
\item La extensión de Raynaud $G^\natural$ sobre $S$ asociada al esquema
  abeliano $A_K$ con reducción semiestable. Dualidad de los esquemas
  abelianos $B,B'$ sobre $S$
\item La extensión de Raynaud $A_K^{\natural\,o}$ en el caso semiestable,
  y el ind-grupo $A_K^\natural$
\item Definición del emparejamiento de monodromía
  $\underline M_\ell\otimes\underline M'_\ell\longrightarrow
  \mathbb Z_{\ell S}$
\item Propiedades del emparejamiento de monodromía. Teorema de
  integralidad y positividad
\item Componentes conexas del modelo de Néron: relación de dualidad,
  comportamiento asintótico
\item Modelo de Néron de una jacobiana y fórmula de Picard--Lefschetz
\item Relaciones con la teoría trascendente: el caso analítico complejo
  \begin{enumerate}[label=\arabic{enumi}.\arabic*.]
  \item Repaso de los grupos analíticos sobre $S$
  \item Completado formal a lo largo de una fibra
  \item La extensión analítica de Raynaud
  \item Interpretación mediante las «estructuras de Hodge mixtas» de
    Deligne
  \item El caso relativamente algebraico
  \end{enumerate}
\item Relaciones con la teoría trascendente: el caso analítico rígido.
  El homomorfismo canónico
  $\underline M_K\longrightarrow A_K^{\natural\,o}$
\end{enumerate}

